% 第7章 自旋表示
\chapter{自旋表示}

\section{Holstein--Primakoff 玻色子}

自旋是矢量算符。在破缺对称相中，其至少一个分量的期望值非零。用自旋绕各自期望值的小涨落来描述有序相是很自然的——这正是自旋波理论的内容。为完成这一任务，Holstein 与 Primakoff 引入玻色子算符 $b$，把三个自旋分量算符表示为

\begin{align*}
S ^ {+} & = \left( \sqrt {2 S - n _ {b}} \, \right) b ,\\
S ^ {-} & = b ^ {\dagger} \sqrt {2 S - n _ {b}} ,\\
S ^ {z} & = - n _ {b} + S ,
\tag{7.1}
\end{align*}

其中 $n_{b} = b^{\dagger}b$。利用 $[b, b^{\dagger}] = 1$，(7.1) 中的算符确实满足自旋对易关系

\begin{equation*}
[ S ^ {\alpha}, S ^ {\beta} ] = i \epsilon^ {\alpha \beta \gamma} S ^ {\gamma} ,
\tag{7.2}
\end{equation*}

拉丁指标取 $x, y, z$，$\epsilon^{\alpha\beta\gamma}$ 是全反对称张量。$b$ 的 Fock 空间太大，物理子空间由态

\begin{equation*}
\{ | n _ {b} \rangle \} _ {S} = \{ | 0 \rangle , | 1 \rangle \dots | 2 S \rangle \}
\tag{7.3}
\end{equation*}

张成。$n_{b} > 2S$ 的非物理 Fock 态由投影算符 $P_{S}$ 消去。在这个子空间中可以验证 $\mathbf{S}^{2} = S(S + 1)$；自旋算符 (7.1) 不会把物理子空间与非物理子空间相连。

Holstein--Primakoff 表示对描述量子 Heisenberg 模型的破缺对称相很有用。(7.1) 中的平方根因子可按 $1/S$ 幂次展开：

\begin{equation*}
\sqrt {2 S - n _ {b}} = \sqrt {2 S} \left( 1 - \frac {n _ {b}}{4 S} - \frac {n _ {b} ^ {2}}{32 S ^ {2}} \dots \right) ,
\tag{7.4}
\end{equation*}

这是自旋涨落绕 $z$ 方向的半经典展开。把 (7.4) 代入哈密顿量并保留到二次项，得到无相互作用的“自旋波”哈密顿量；更高阶项引入自旋波之间的相互作用。截断会把物理子空间与非物理子空间耦合起来。尽管如此，低阶自旋波近似已被证明在 Heisenberg 铁磁体与反铁磁体的某些有序相中极为成功。文献注释中列出这一主题的若干里程碑论文。第 11 章将以两种方式导出自旋波理论：从自旋相干态路径积分，以及展开 Holstein--Primakoff 算符。

\section{Schwinger 玻色子}

Heisenberg 模型的对称相，用能显式展现哈密顿量转动不变性的表示来描述更容易。两个 Schwinger 玻色子 $a$、$b$ 把自旋算符表示为

\begin{align*}
S ^ {x} + i S ^ {y} & = a ^ {\dagger} b ,\\
S ^ {x} - i S ^ {y} & = b ^ {\dagger} a ,\\
S ^ {z} & = \frac {1}{2} \left( a ^ {\dagger} a - b ^ {\dagger} b \right) .
\tag{7.5}
\end{align*}

容易验证如上定义的自旋分量满足 (7.2)。自旋大小 $S$ 定义物理子空间

\begin{equation*}
\{ | n _ {a}, n _ {b} \rangle : n _ {a} + n _ {b} = 2 S \}.
\tag{7.6}
\end{equation*}

该子空间由投影算符（或“约束”）$P_{S}$ 给出：

\begin{equation*}
P _ {S} \left( a ^ {\dagger} a + b ^ {\dagger} b - 2 S \right) = 0.
\tag{7.7}
\end{equation*}

在投影后的子空间中，自旋大小是确定的：

\begin{equation*}
\mathbf {S} ^ {2} P _ {S} = S ( S + 1 ) P _ {S}.
\tag{7.8}
\end{equation*}

图~\ref{fig:7.1} 在 Fock 空间 $n_a, n_b$ 中画出了 $P_S$ 投影的子空间。

自旋态由

\begin{equation*}
\left| S, m \right\rangle = \frac {\left( a ^ {\dagger} \right) ^ {S + m}}{\sqrt {( S + m ) !}} \, \frac {\left( b ^ {\dagger} \right) ^ {S - m}}{\sqrt {( S - m ) !}} \left| 0 \right\rangle ,
\tag{7.9}
\end{equation*}

给出，其中 $|S,m\rangle$ 标记 $S^{2}$ 与 $S^{z}$ 的本征值，$|0\rangle$ 是 Schwinger 玻色子真空。例如，自旋 $\frac{1}{2}$ 的态在二次量子化记号下为

\begin{align*}
\left| \uparrow \right\rangle & = a ^ {\dagger} \left| 0 \right\rangle ,\\
\left| \downarrow \right\rangle & = b ^ {\dagger} \left| 0 \right\rangle .
\tag{7.10}
\end{align*}

\begin{figure}[H]
\centering
\includegraphics[width=0.5\textwidth]{fig7_1.jpg}
\caption{Schwinger 玻色子 Fock 空间中自旋 $S$ 的投影子空间。}
\label{fig:7.1}
\end{figure}

Schwinger 玻色子便于计算自旋算符的矩阵元：由于不含平方根（与 Holstein--Primakoff 玻色子不同），自旋算符在 Fock 态之间的矩阵元像自由玻色子一样因子化。然而这并不必然简化非 Fock 波函数中自旋关联的计算：Hilbert 空间上的局域约束 (7.6) 引入了 $a$、$b$ 占据数之间的关联。

把表示 (7.5) 推广到 $N$ 个“味”，就把 $SU(N)$ 的生成元定义为广义“自旋”，见第 16 章。大 $N$ 推广允许由小参数 $1/N$ 控制的简单平均场理论，将在第 17、18 章综述。大 $N$ 平均场理论由实际上无相互作用的 Bose 准粒子描述，不同味之间的关联被忽略。

Schwinger 玻色子（SB）与 Holstein--Primakoff 玻色子 (7.1)（HP）密切相关。利用约束 (7.7) 消去 $a$ 玻色子，得到对应关系

\begin{equation*}
\begin{array}{ccc}
\text{SB} & & \text{HP}\\
b & \leftrightarrow & b ,\\
a & \leftrightarrow & \sqrt {2 S - n _ {b}} .
\end{array}
\tag{7.11}
\end{equation*}

Schwinger 玻色子在自旋空间中提供对称表示，而 Holstein--Primakoff 玻色子挑出 $S^{z}$ 方向（后者定义其真空）。因此两种表示适合不同的近似方案：HP 用于破缺对称相，SB 用于对称相。

\subsection{自旋转动}

自旋算符是 SU(2) 变换群的生成元。$\mathcal{R}$ 由三个 Euler 角 $\phi, \theta, \chi$ 参数化：

\begin{equation*}
\mathcal {R} = e ^ {i \phi S ^ {z}} e ^ {i \theta S ^ {y}} e ^ {i \chi S ^ {z}}.
\tag{7.12}
\end{equation*}

自旋算符是正规双线性算符（见 A.2 节）。Schwinger 玻色子产生算符作为 SU(2) 中的矢量变换\footnote{见附录 A 的 (A.17)。}：

\begin{align*}
\binom {a ^ {\dagger}} {b ^ {\dagger}} ' & = \mathcal {R} \binom {a ^ {\dagger}} {b ^ {\dagger}} \mathcal {R} ^ {- 1}\\
& = e ^ {i \frac {1}{2} \chi \sigma^ {z}} e ^ {i \frac {1}{2} \theta \sigma^ {y}} e ^ {i \frac {1}{2} \phi \sigma^ {z}} \binom {a ^ {\dagger}} {b ^ {\dagger}}\\
& = \left( \begin{array}{cc} u \, e ^ {i \chi / 2} & v \, e ^ {i \chi / 2} \\ - v ^ {*} e ^ {- i \chi / 2} & u ^ {*} e ^ {- i \chi / 2} \end{array} \right) \binom {a ^ {\dagger}} {b ^ {\dagger}} ,
\tag{7.13}
\end{align*}

其中函数 $u$、$v$ 定义为

\begin{align*}
u ( \theta , \phi ) & = \cos ( \theta / 2 ) \, e ^ {i \phi / 2} ,\\
v ( \theta , \phi ) & = \sin ( \theta / 2 ) \, e ^ {- i \phi / 2} .
\tag{7.14}
\end{align*}

\section{自旋相干态}

自旋相干态是把 (7.12) 的转动算符 $\mathcal{R}$ 作用于最大极化态 $|S, S\rangle$ 得到的一族自旋态：

\begin{align*}
\left| \hat {\Omega} \right\rangle & = \mathcal {R} ( \chi , \theta , \phi ) \left| S, S \right\rangle\\
& = e ^ {i S ^ {z} \phi} e ^ {i S ^ {y} \theta} e ^ {i S ^ {z} \chi} \left| S, S \right\rangle ,
\tag{7.15}
\end{align*}

其中单位矢量

\begin{equation*}
\hat {\Omega} = \left( \sin \theta \cos \phi \, , \; \sin \theta \sin \phi \, , \; \cos \theta \right)
\tag{7.16}
\end{equation*}

参数化该自旋相干态。我们可以任意定义 $\chi$——这是一种规范自由，可以通过固定 $\chi$ 消去。两个独立角度是“纬度”$\theta$ 与“经度”$\phi$，定义域为

\begin{equation*}
\{ \theta \in [ 0, \pi ] \, , \; \phi \in [ - \pi , \pi ) \}.
\tag{7.17}
\end{equation*}

利用 (7.13)，Schwinger 玻色子变换到转动坐标系 $a \rightarrow a_{\theta,\phi}'$，给出相干态的显式表示：

\begin{align*}
\left| \hat {\Omega} \right\rangle & = e ^ {i S \chi} \, \frac {\left( a ^ {\dagger '} \right) ^ {2 S}}{\sqrt {( 2 S ) !}} \left| 0 \right\rangle = e ^ {i S \chi} \, \frac {\left( u \, a ^ {\dagger} + v \, b ^ {\dagger} \right) ^ {2 S}}{\sqrt {( 2 S ) !}} \left| 0 \right\rangle\\
& = e ^ {i S \chi} \sqrt {( 2 S ) !} \, \sum_ {m} \frac {u ^ {S + m} v ^ {S - m}}{\sqrt {( S + m ) ! \, ( S - m ) !}} \left| S, m \right\rangle ,
\tag{7.18}
\end{align*}

其中 $u$、$v$ 已在 (7.14) 定义。

利用 (7.18) 可以计算任意两个相干态的交叠：

\begin{align*}
\langle \hat {\Omega} | \hat {\Omega} ' \rangle & = e ^ {- i S ( \chi - \chi ' )} \left( u ^ {*} u ' + v ^ {*} v ' \right) ^ {2 S}\\
& = \left( \frac {1 + \hat {\Omega} \cdot \hat {\Omega} '}{2} \right) ^ {S} e ^ {- i S \psi}\\
\psi & = 2 \arctan \left[ \tan \left( \frac {\phi - \phi '}{2} \right) \frac {\cos \left[ \frac {1}{2} ( \theta + \theta ' ) \right]}{\cos \left[ \frac {1}{2} ( \theta - \theta ' ) \right]} \right] + \chi - \chi ' ,
\tag{7.19}
\end{align*}

其中 $\chi, \chi'$ 依赖规范约定。

如 (7.19) 所示，相干态并不正交。现在证明它们张成自旋 $S$ 的态空间。群参数上的积分测度定义为

\begin{equation*}
\frac {2 S + 1}{4 \pi} d \hat {\Omega} = \frac {2 S + 1}{4 \pi} d \theta \sin \theta \, d \phi .
\tag{7.20}
\end{equation*}

这是 SU(2) Lie 群的 Haar 测度。利用这一测度，单位分解由积分\footnote{{$\theta$} 积分的计算见 7.3.1 小节。}给出

\begin{align*}
& \frac {2 S + 1}{4 \pi} \int d \hat {\Omega} \left| \hat {\Omega} \right\rangle \left\langle \hat {\Omega} \right|\\
& = \frac {( 2 S + 1 )}{2} \int_ {- 1} ^ {1} d \cos \theta \, \sum_ {m} \left( \frac {1 + \cos \theta}{2} \right) ^ {S + m} \left( \frac {1 - \cos \theta}{2} \right) ^ {S - m}\\
& \quad \times \frac {( 2 S ) !}{( S + m ) ! ( S - m ) !} \left| S, m \right\rangle \left\langle S, m \right|\\
& = \sum_ {m} \left| S, m \right\rangle \left\langle S, m \right| = I .
\tag{7.21}
\end{align*}

因此相干态构成一个超完备基底。

用同样方式还可证明另一个有用关系\footnote{见习题 3。}

\begin{equation*}
\frac {( S + 1 ) ( 2 S + 1 )}{4 \pi} \int d \hat {\Omega} \, \hat {\Omega} ^ {\alpha} \left| \hat {\Omega} \right\rangle \left\langle \hat {\Omega} \right| = \mathbf {S} ^ {\alpha} \, , \quad \alpha = x, y, z.
\tag{7.22}
\end{equation*}

上述关系容易推广到多自旋。考虑 $\mathcal{N}$ 个格点上自旋的格子，格点记作 $i$。多自旋相干态是单自旋态的乘积：

\begin{equation*}
\left| \hat {\Omega} \right\rangle = \prod_ {i = 1} ^ {\mathcal {N}} \left| \hat {\Omega} _ {i} \right\rangle .
\tag{7.23}
\end{equation*}

其交叠为

\begin{equation*}
\langle \hat {\Omega} | \hat {\Omega} ' \rangle = \prod_ {i} \left( \frac {1 + \hat {\Omega} _ {i} \cdot \hat {\Omega} _ {i} '}{2} \right) ^ {S} e ^ {- i S \sum_ {i} \psi [ \hat {\Omega} _ {i}, \hat {\Omega} _ {i} ' ]}.
\tag{7.24}
\end{equation*}

乘积空间中的单位分解为

\begin{equation*}
\int \prod_ {i} \left( \frac {2 S + 1}{4 \pi} d \hat {\Omega} _ {i} \right) \left| \hat {\Omega} \right\rangle \left\langle \hat {\Omega} \right| = I.
\tag{7.25}
\end{equation*}

任何波函数 $\Psi$ 的自旋关联都可以用积分计算：

\begin{align*}
\frac {\langle \Psi | \mathbf {S} _ {i} \cdot \mathbf {S} _ {j} | \Psi \rangle}{\langle \Psi | \Psi \rangle} & = \frac {( S + 1 - \delta_ {ij} ) ( S + 1 )}{Z} \int \prod_ {i} d \hat {\Omega} _ {i} \left| \Psi [ \hat {\Omega} ] \right| ^ {2} \hat {\Omega} _ {i} \cdot \hat {\Omega} _ {j} ,\\
Z & = \int \prod_ {i} d \hat {\Omega} _ {i} \left| \Psi [ \hat {\Omega} ] \right| ^ {2} ,
\tag{7.26}
\end{align*}

这里用到了 (7.22)。表示 $\Psi[\hat{\Omega}]=\langle\Psi|\hat{\Omega}\rangle$ 是连续的（像 Schrödinger 波函数一样），习题 4 将证明所有自旋算符都表示为变量 $u(\theta,\phi)$、$v(\theta,\phi)$ 的微分算符。自旋相干态可用于计算任何算符的迹。若 $|n\rangle$ 是任意正交归一基，则

\begin{align*}
\operatorname{Tr} \mathcal {O} & = \sum_ {n} \langle n | \mathcal {O} | n \rangle\\
& = \left( \frac {2 S + 1}{4 \pi} \right) ^ {2} \int d \hat {\Omega} \int d \hat {\Omega} ' \sum_ {n} \langle n | \hat {\Omega} \rangle \langle \hat {\Omega} | \mathcal {O} | \hat {\Omega} ' \rangle \langle \hat {\Omega} ' | n \rangle\\
& = \frac {2 S + 1}{4 \pi} \int d \hat {\Omega} \, \langle \hat {\Omega} | \mathcal {O} | \hat {\Omega} \rangle .
\tag{7.27}
\end{align*}

相干态阐明了经典自旋与量子自旋之间的对应。经典极限由 $S \rightarrow \infty$ 得到。在该极限下，由 (7.24)，不同相干态的交叠随 $S$ 指数消失；自旋算符的期望值成为单位矢量的函数，与经典自旋毫无二致。量子效应因此与相干态的非正交性相联系——这也意味着 $\Psi[\hat{\Omega}]$ 在 $\hat{\Omega}$ 空间中有有限宽度。

\subsection{$\theta$ 积分}

把 (7.21) 中的 $\theta$ 积分定义为

\begin{equation*}
I _ {S, m} = \frac {1}{2} \int_ {0} ^ {\pi} d \theta \sin \theta \left( \frac {1 + \cos \theta}{2} \right) ^ {S + m} \left( \frac {1 - \cos \theta}{2} \right) ^ {S - m}.
\tag{7.28}
\end{equation*}

换积分变量

\begin{equation*}
x = \frac {\cos \theta + 1}{2} ,
\tag{7.29}
\end{equation*}

可以证明

\begin{align*}
I _ {S, m} & = \int_ {0} ^ {1} d x \, x ^ {S + m} ( 1 - x ) ^ {S - m}\\
& = ( S + m )! \, ( S - m )! \, / \, ( 2 S + 1 )! .
\tag{7.30}
\end{align*}

验证方法是构造生成函数

\begin{equation*}
f _ {S} ( z ) = \sum_ {n = 0} ^ {2 S} \frac {( 2 S ) !}{( 2 S - n )! \, n !} \, I _ {S, n - S} \, z ^ {n}.
\tag{7.31}
\end{equation*}

把 (7.30) 代入 (7.31) 并用二项式展开：

\begin{align*}
( 2 S + 1 ) f _ {S} ( z ) & = ( z ^ {2 S + 1} - 1 ) / ( z - 1 )\\
& = 1 + z + \ldots + z ^ {2 S} .
\tag{7.32}
\end{align*}

由 (7.31)，$I_{S,m}$ 正是 $f_{S}(z)$ 中 $z^{m + S}$ 的系数，这就证明了 (7.30)。

\begin{exercises}

\item 直接从定义 (7.15) 验证 $|\hat{\Omega}\rangle$ 是 $\hat{\Omega}$ 方向自旋分量的本征态：

\begin{equation*}
\hat {\Omega} \cdot \mathbf {S} \left| \hat {\Omega} \right\rangle = S \left| \hat {\Omega} \right\rangle .
\tag{7.33}
\end{equation*}

提示：利用三个自旋算符 $(S^{x}, S^{y}, S^{z})$ 的 $O(3)$（矢量）变换性质。

\item 利用 (7.33)，证明 Heisenberg 模型在相干态中的期望值为

\begin{equation*}
H [ \hat {\Omega} ] = \langle \hat {\Omega} | \mathcal {H} | \hat {\Omega} \rangle = \frac {S ^ {2}}{2} \sum_ {i, j} J _ {ij} \, \hat {\Omega} _ {i} \cdot \hat {\Omega} _ {j}.
\tag{7.34}
\end{equation*}

提示：对每个自旋，把它的分量写成 $\hat{\Omega}$ 方向的分量与两个横向分量。

\item 循 (7.21) 的推导，证明恒等式 (7.22)。

\item 证明对单自旋的任意波函数，下列关系成立：

\begin{align*}
\langle \Psi | S ^ {+} | \hat {\Omega} \rangle & = v \, \partial_ {u} \Psi ( \hat {\Omega} ) ,\\
\langle \Psi | S ^ {-} | \hat {\Omega} \rangle & = u \, \partial_ {v} \langle \Psi | \hat {\Omega} \rangle ,\\
\langle \Psi | S ^ {z} | \hat {\Omega} \rangle & = \frac {1}{2} \left( u \, \partial_ {u} - v \, \partial_ {v} \right) \Psi ( \hat {\Omega} ) .
\tag{7.35}
\end{align*}

\item 给定微分算符

\begin{equation*}
\mathcal {T} = \sum_ {klj} T _ {klj} \, \partial_ {u} ^ {k} \partial_ {v} ^ {l} u ^ {k + j} v ^ {l - j} ,
\tag{7.36}
\end{equation*}

它在任意两个自旋 $S$ 波函数之间的矩阵元为

\begin{equation*}
\langle \Psi | \mathcal {T} | \Psi ' \rangle = \frac {2 S + 1}{4 \pi} \int d \hat {\Omega} \, \Psi ^ {*} ( \hat {\Omega} ) \left[ \mathcal {T} _ {u, v} \Psi ' ( \hat {\Omega} ) \right].
\tag{7.37}
\end{equation*}

证明可以通过代入

\begin{equation*}
\langle \Psi | \mathcal {T} | \Psi ' \rangle = \frac {2 S + 1}{4 \pi} \int d \hat {\Omega} \, \Psi ^ {*} ( \hat {\Omega} ) \, T [ u ^ {*}, v ^ {*}, u, v ] \, \Psi ' ( \hat {\Omega} ) ,
\tag{7.38}
\end{equation*}

来计算 (7.37)，其中

\begin{equation*}
T \left[ u ^ {*}, v ^ {*}, u, v \right] = \sum_ {klj} C _ {kl} \, T _ {klj} \, u ^ {*k} v ^ {*l} u ^ {k + j} v ^ {l - j}
\tag{7.39}
\end{equation*}

而

\begin{equation*}
C _ {kl} = \prod_ {m = 2} ^ {k + l + 1} ( 2 S + m ).
\tag{7.40}
\end{equation*}

提示：把 $\Psi$、$\Psi'$ 展开成 $u, v$ 的二项式，对 (7.37) 与 (7.40) 先积 $d\phi$ 再积 $d\theta$。

\end{exercises}

\begin{chapbib}
\item Holstein--Primakoff 算符定义于
  T. Holstein and H. Primakoff, Phys. Rev. \textbf{58}, 1908 (1940)。
\item Schwinger 玻色子的定义及其在 Clebsch--Gordan 系数计算中的应用：
  J. Schwinger, \textit{Quantum Theory of Angular Momentum}, edited by L. Biedenharn and H. Van Dam (Academic, New York, 1965)。
\item 自旋相干态及其推广见
  A. Perelomov, \textit{Generalized Coherent States and Their Applications} (Springer-Verlag, New York, 1986)；
  J. R. Klauder and B. Skagerstam, \textit{Coherent States} (World Scientific, 1985)。
\item 自旋相干态与半经典近似的联系见
  J. Klauder, Phys. Rev. D \textbf{19}, 2349 (1979)。
\item 相干态在量子 Heisenberg 铁磁体中的应用见
  F. D. M. Haldane, Phys. Rev. Lett. \textbf{57}, 1488 (1986)。
\end{chapbib}
